% Generated draft: requires mathematical and source review.
\documentclass[11pt]{article}
\usepackage[margin=25mm]{geometry}
\usepackage{fontspec}
\setmainfont{Latin Modern Roman}
\usepackage{amsmath,amssymb,mathtools,unicode-math}
\setmathfont{Latin Modern Math}
\usepackage{enumitem,needspace,xurl,hyperref}
\setlength{\parindent}{0pt}
\setlength{\parskip}{5pt}
\setlength{\emergencystretch}{3em}
\newcommand{\N}{\mathbb{N}}
\newcommand{\Z}{\mathbb{Z}}
\newcommand{\Q}{\mathbb{Q}}
\newcommand{\R}{\mathbb{R}}
\newcommand{\C}{\mathbb{C}}
\newcommand{\F}{\mathbb{F}}
\newcommand{\A}{\mathbb{A}}
\newcommand{\PP}{\mathbb{P}}
\newcommand{\E}{\mathbb{E}}
\newcommand{\eps}{\varepsilon}
\newcommand{\dd}{\,\mathrm{d}}
\begin{document}
\section*{UnsolvedMath ID: 6200063}
\textbf{Source problem code:} catalogue entry
\subsection{Definitions and mathematical statement}
Let $G$ be a word-hyperbolic group and let $(X,d,\mu)$ be the intended metric-measure realization of $X=\partial_\infty G$. In a usual $Q$-Loewner setting, $Q>1$, the measure is Ahlfors $Q$-regular and the $Q$-modulus of curves joining disjoint nondegenerate continua has a positive lower bound depending only on their relative distance. The source does not fix $Q$, the metric, or the measure.

A nonnegative Borel function $g$ is an upper gradient of a real function $u$ if every rectifiable curve $\gamma$ joining $x$ to $y$ satisfies
\[
 |u(x)-u(y)|\le\int_\gamma g\,\mathrm ds.
\]
Under the usual weak $(1,1)$-Poincaré convention, supporting a $1$-Poincaré inequality means that constants $C>0$ and $\lambda\ge1$ exist so that, for every ball $B=B(x,r)$, admissible function $u$, and its upper gradient $g$,
\[
 \frac1{\mu(B)}\int_B|u-u_B|\,\mathrm d\mu
 \le Cr\,\frac1{\mu(\lambda B)}\int_{\lambda B}g\,\mathrm d\mu,
 \qquad
 u_B=\frac1{\mu(B)}\int_Bu\,\mathrm d\mu.
\]
One can initially test this with Lipschitz $u$; the full test class, weak-upper-gradient variant, and whether dilation is allowed require specification.

Heinonen's question is to determine assumptions on $G$, beyond having a $Q$-Loewner boundary, that guarantee the boundary supports the intended $1$-Poincaré inequality. The exponent $1$ concerns the upper-gradient integral; a Poincaré inequality with a different exponent does not settle this endpoint question. No particular extra group hypothesis or necessary-and-sufficient characterization is supplied in the source.
\subsection{Short English statement}
Which additional properties of a hyperbolic group with a $Q$-Loewner boundary force an endpoint $1$-Poincaré inequality on that boundary? Metric, measure, and the precise Poincaré convention must be fixed.

The supplied summary mentions Bourdon--Kleiner and Poincaré consequences of Loewner properties but gives no exact exponent, theorem, or reference. It calls the $p=1$ question subtle and open. The OPEN-TRIAGE label and 2005 source are retained with current status explicitly marked for review.
\subsection{Sources}
\begin{itemize}
\item Local UnsolvedMath registry, permanent ID 6200063, source code catalogue entry, ``Boundaries of Groups and Kleinian Groups --- Problem 63,'' proposed by Juha Heinonen, 2005. Dataset curation and metadata: Ulam AI and UnsolvedMath Contributors, CC BY 4.0.
\item Kapovich, \emph{Problems on Boundaries of Groups and Kleinian Groups} (2005), Problem 63, PDF page 17, \url{https://www.math.ucdavis.edu/~kapovich/EPR/problems.pdf}; normalized pdf-text extraction, accessed 2026-07-29.
\item The supplied research summary mentions Bourdon--Kleiner without complete bibliographic details or a precise endpoint result. Source-specific provenance marks current status and statement-reuse terms NEEDS\_REVIEW for the public author-hosted PDF, separately from the dataset licence.
\end{itemize}
\end{document}
